\section{From a Model to an Agent}
\paragraph{Purpose}
A language model only produces text. An \textbf{agent} is an LLM placed \tB{in a loop} with
\tB{tools} (search, code execution, files, APIs) that acts on an environment over several
steps to reach a goal, using the results of its actions to decide the next one.\\
Distinction: in a \textbf{workflow} the steps are fixed by the code (the LLM fills each
step), in an \textbf{agent} the LLM itself \uB{decides the steps and the tools}. The agent is
more flexible but slower, more expensive and less predictable, so the simplest solution
(a single call, possibly with retrieval) should be tried first.

\section{Tool Use}
\paragraph{Theory}
Each tool is described to the model by a \emph{name}, a \emph{description} and a
\emph{schema} of its parameters (JSON), placed in the prompt.
\begin{enumerate}
    \item the model decides to call a tool and outputs a \tB{structured call}: the name and
        the arguments;
    \item the \uB{application, not the model}, validates and executes it;
    \item the result is appended to the conversation, and the model continues.
\end{enumerate}
The model never executes anything: this is the point where permissions, validation and
sandboxing are enforced. The behavior is learned by fine-tuning on examples of calls
(Toolformer trained a model to insert API calls in text by itself), and standards such as
the Model Context Protocol (MCP) describe how tools and data sources are connected to
models.

\paragraph{Weaknesses} The model may call a tool that does not exist, give invalid or
invented arguments, or misread the result: calls must be validated, and errors returned to
the model so that it can correct itself.

\section{The Agent Loop: ReAct}
\paragraph{Theory}
ReAct (Yao et al., 2022) interleaves \emph{reasoning} and \emph{acting}: the model writes a
thought, an action, receives an observation, and repeats.
\begin{center}
\fbox{\parbox{.9\textwidth}{
\texttt{context} $\leftarrow$ goal + tool descriptions\\
\textbf{repeat} until a final answer or a maximum number of steps:\\
\hspace*{1em}\texttt{output} $\leftarrow$ LLM(\texttt{context})\\
\hspace*{1em}\textbf{if} \texttt{output} is a tool call: \texttt{result} $\leftarrow$ execute it;
append \texttt{output} and \texttt{result} to \texttt{context}\\
\hspace*{1em}\textbf{else}: return \texttt{output}
}}
\end{center}

\paragraph{Reinforcement learning view}
This is the setting of the RL chapter: the \emph{state} is the history in the context, the
\emph{action} is a text or a tool call, the \emph{observation} is the tool output, and the
\emph{policy} is the LLM. The state is only partially observed (the world is not entirely in
the context), and the reward arrives at the end (the task is solved or not). Training the
model with verifiable rewards on such tasks is an application of the GRPO or PPO methods.

\section{Planning, Reflection and Multiple Agents}
\begin{itemize}
    \item \textbf{planning}: decompose the goal in sub-tasks, either first (plan then
        execute) or step by step, and revise the plan from the observations;
    \item \textbf{reflection}: the model critiques its own result or the error message, and
        tries again;
    \item \textbf{orchestrator and workers}: a main agent delegates sub-tasks to other
        calls, each one with its own short context, and collects their results.
\end{itemize}

\section{Memory and Context}
\begin{itemize}
    \item \textbf{short-term memory}: the context window, limited to $n_{ctx}$ tokens. Every
        step sends the whole history again, and the \uB{KV cache} of the common prefix avoids
        recomputing it;
    \item \textbf{long-term memory}: notes or documents stored outside and retrieved when
        needed (see RAG);
    \item when the window fills, the old steps are \emph{summarized} or the large tool
        outputs are removed: a bad summary loses information the agent will need.
\end{itemize}

\section{Reliability and Evaluation}
\paragraph{Errors compound}
If each step succeeds with probability $p$ independently, a task of $n$ steps succeeds with
probability about $p^{n}$: for $p=0.95$ and $n=20$ it is $0.95^{20}\approx0.36$. Long
horizons need reliable steps, checks and recovery.

\paragraph{Metrics}
\begin{itemize}
    \item the \textbf{success rate} on tasks with an automatic check (tests that pass, the
        final state of the environment);
    \item $pass@k$: the probability that at least one of $k$ attempts succeeds. From $n$
        attempts of which $c$ succeed, the unbiased estimate is
        $1-\binom{n-c}{k}\Big/\binom{n}{k}$;
    \item the \textbf{cost}: tokens, number of steps and tool calls, latency.
\end{itemize}
Typical failure modes are loops without progress, a wrong tool or arguments, a hallucinated
result, and an error early in the run that propagates.

\section{Safety}
\begin{itemize}
    \item \textbf{indirect prompt injection}: instructions hidden in a web page, a file or a
        tool output are read by the model as if they were legitimate, and can make it take
        actions the user did not want. The text of tool results must be treated as
        \tB{data, not as instructions};
    \item \textbf{least privilege}: give the agent only the tools and the permissions the
        task needs, and run the code in a sandbox;
    \item \textbf{human in the loop} before irreversible or external actions (sending, paying,
        deleting);
    \item limits on the number of steps and on the budget, and logs of every call to audit
        what the agent did.
\end{itemize}
